% =====================================================================
% Capítulo 4 — Esboço de seção/interlúdio (análise de rede textual)
% Modelo: esboco_secao_infranodus.tex (capítulo 1).
%
% Estilo seguido (regras da tese): primeira pessoa do singular, sem
% travessões em texto principal, "capítulo" minúsculo em referências
% cruzadas, aspas com \enquote{}, sem fórmulas do tipo "não X, mas Y".
%
% Observação técnica: copie
%   infranodus/cap4/infranodus_cap4_focus.png
% para
%   figuras/cap.4/infranodus_cap4_focus.png
% =====================================================================

\subsection*{Interlúdio: a rede textual do capítulo}

Submeto por fim à análise de rede textual do capítulo~\ref{capitulo1} os
32.745~\textit{tokens} deste capítulo, o mais extenso da tese. O
grafo organiza-se em torno de \textit{spira}, que é ao mesmo tempo o
termo de maior \textit{degree} e, com larga margem, a maior ponte de
todos os capítulos (\textit{betweenness}~0,55). Ao seu redor
condensa-se o aparato latouriano da inscrição, e os pares de maior
associação são exatamente os conceitos-chave que mobilizo:
\textit{caixa}~$\leftrightarrow$~\textit{preta}, a caixa-preta, e
\textit{imutável}~$\leftrightarrow$~\textit{móvel} (NPMI~0,84), o
móvel imutável de \textcite{Latour1986Visualization}. O par
\textit{cálculo}~$\leftrightarrow$~\textit{centro} (0,77) inscreve
na rede o \enquote{centro de cálculo}, e a cena clínica aparece na
co-ocorrência mais pesada do capítulo,
\textit{respiratória}~$\leftrightarrow$~\textit{insuficiência}
(peso~308, NPMI~0,84).

A cadeia de translações que o capítulo narra deixa-se ler na sequência
de comunidades do grafo. Ela parte da cena clínica (\textit{respiratória},
\textit{insuficiência}, \textit{covid}, \textit{paciente},
\textit{ruído}), passa pelo sinal (\textit{espectrograma},
\textit{sinal}, \textit{acústico}), atravessa o dado e o modelo
(\textit{dado}, \textit{modelo}, \textit{áudio},
\textit{rede}~$\leftrightarrow$~\textit{neural}) e chega ao
\textit{artigo}. As pontes do grafo, além de \textit{spira}, são
\textit{inscrição}, \textit{covideiro}, \textit{cadeia}, \textit{artigo}
e \textit{objeto}, e vale deter-me em \textit{covideiro}: o neologismo
que cunho no capítulo, o actante pandêmico, opera como conector de
comunidades, de sorte que a própria invenção lexical do texto faz, na
rede, trabalho de tradução. A teoria e o objeto empírico, que nos
capítulos anteriores custavam a se tocar, aqui compartilham vizinhança,
e o grafo mostra a inscrição latouriana já colada à matéria acústica e
estatística do SPIRA.

\begin{figure}[!htb]
\centering
\includegraphics[width=1.0\textwidth,keepaspectratio]{figuras/cap.4/infranodus_cap4_focus.png}
\caption[Rede textual do capítulo~\ref{capitulo4} (\textit{text network
analysis})]{Rede textual do capítulo~\ref{capitulo4}
(\textit{text network analysis}). Os nós são os termos de maior
\textit{degree} ponderado; o tamanho do nó acompanha o \textit{degree};
a cor indica a comunidade detectada por Louvain ponderado; as arestas,
arqueadas, recebem a mistura das cores das comunidades que ligam.
Versão de alta resolução, arquivo \texttt{.gexf} para o \textit{Gephi} e
relatório completo em \texttt{infranodus/cap4/} no repositório desta
tese.}
\label{fig:infranodus-cap4}
\end{figure}

As ligações mais fracas caem, de modo revelador, nas juntas da própria
cadeia. A lacuna mais forte separa a teoria da inscrição
(\textit{inscrição}, \textit{cadeia}, \textit{dispositivo}) do par
\textit{dado}~$\leftrightarrow$~\textit{modelo}; e seguem ralas as
costuras entre o espectrograma e o artigo, e entre o sinal e o dado. A
rede é mais rala exatamente onde o capítulo promete o essencial, que é
mostrar como cada elo converte um estado no seguinte, da voz ao sinal, do
sinal ao espectrograma, do espectrograma ao dado, do dado ao modelo e do
modelo ao artigo. Leio essa rarefação como a agenda de reescrita que o
capítulo me deixa, tornando explícita cada tradução em vez de confiá-la à
adjacência dos termos. O gesto latouriano me dá o critério, pois uma
cadeia de inscrições só se sustenta quando cada passagem é reversível e
narrável \parencite{Latour1986Visualization}, e o grafo assinala os
pontos em que essa narração ainda me falta.
